% Options for packages loaded elsewhere
\PassOptionsToPackage{unicode}{hyperref}
\PassOptionsToPackage{hyphens}{url}
\documentclass[
  ignorenonframetext,
]{beamer}
\newif\ifbibliography
\usepackage{pgfpages}
\setbeamertemplate{caption}[numbered]
\setbeamertemplate{caption label separator}{: }
\setbeamercolor{caption name}{fg=normal text.fg}
\beamertemplatenavigationsymbolsempty
% remove section numbering
\setbeamertemplate{part page}{
  \centering
  \begin{beamercolorbox}[sep=16pt,center]{part title}
    \usebeamerfont{part title}\insertpart\par
  \end{beamercolorbox}
}
\setbeamertemplate{section page}{
  \centering
  \begin{beamercolorbox}[sep=12pt,center]{section title}
    \usebeamerfont{section title}\insertsection\par
  \end{beamercolorbox}
}
\setbeamertemplate{subsection page}{
  \centering
  \begin{beamercolorbox}[sep=8pt,center]{subsection title}
    \usebeamerfont{subsection title}\insertsubsection\par
  \end{beamercolorbox}
}
% Prevent slide breaks in the middle of a paragraph
\widowpenalties 1 10000
\raggedbottom
\AtBeginPart{
  \frame{\partpage}
}
\AtBeginSection{
  \ifbibliography
  \else
    \frame{\sectionpage}
  \fi
}
\AtBeginSubsection{
  \frame{\subsectionpage}
}
\usepackage{iftex}
\ifPDFTeX
  \usepackage[T1]{fontenc}
  \usepackage[utf8]{inputenc}
  \usepackage{textcomp} % provide euro and other symbols
\else % if luatex or xetex
  \usepackage{unicode-math} % this also loads fontspec
  \defaultfontfeatures{Scale=MatchLowercase}
  \defaultfontfeatures[\rmfamily]{Ligatures=TeX,Scale=1}
\fi
\usepackage{lmodern}
\ifPDFTeX\else
  % xetex/luatex font selection
\fi
% Use upquote if available, for straight quotes in verbatim environments
\IfFileExists{upquote.sty}{\usepackage{upquote}}{}
\IfFileExists{microtype.sty}{% use microtype if available
  \usepackage[]{microtype}
  \UseMicrotypeSet[protrusion]{basicmath} % disable protrusion for tt fonts
}{}
\makeatletter
\@ifundefined{KOMAClassName}{% if non-KOMA class
  \IfFileExists{parskip.sty}{%
    \usepackage{parskip}
  }{% else
    \setlength{\parindent}{0pt}
    \setlength{\parskip}{6pt plus 2pt minus 1pt}}
}{% if KOMA class
  \KOMAoptions{parskip=half}}
\makeatother
\setlength{\emergencystretch}{3em} % prevent overfull lines
\providecommand{\tightlist}{%
  \setlength{\itemsep}{0pt}\setlength{\parskip}{0pt}}
\usepackage{bookmark}
\IfFileExists{xurl.sty}{\usepackage{xurl}}{} % add URL line breaks if available
\urlstyle{same}
\hypersetup{
  hidelinks,
  pdfcreator={LaTeX via pandoc}}

\author{\texorpdfstring{}{}}
\date{}

\begin{document}

\begin{frame}[fragile]{Abstract}
\protect\phantomsection\label{abstract}
The Docxology Template is a modular, high-integrity research environment
that unifies computational experimentation, academic manuscript
production, and cryptographic provenance verification within a single
reproducible pipeline. Built on a Two-Layer Architecture separating ten
reusable infrastructure subpackages---totaling 137 Python modules and
validated by over 3,000 infrastructure tests---from self-contained
project workspaces, the system orchestrates an eight-stage build
pipeline that progresses from environment sanitization through test
execution, analysis script invocation, Pandoc/XeLaTeX rendering,
SHA-256/512 hashing with alpha-channel steganographic watermarking,
structural PDF validation, and LLM-assisted review to produce
peer-review-ready manuscripts with embedded tamper-evident provenance.
Unlike existing workflow managers (Snakemake, Nextflow, CWL) that focus
on computational pipeline orchestration, and literate programming tools
(Quarto, Jupyter Book, R Markdown) that focus on document rendering, the
Docxology Template integrates both concerns with automated testing,
documentation generation, and cryptographic provenance in a single
coherent system. We enforce a Zero-Mock testing policy where pipeline
advancement is contingent on 90\% project-level and 60\%
infrastructure-level coverage thresholds, using exclusively real
filesystem operations, subprocess invocations, and network calls---no
mock objects, no test doubles. A Documentation Duality standard equips
every directory with both human-readable \texttt{README.md} files and
machine-readable \texttt{AGENTS.md} files optimized for AI coding
assistants, enabling a novel human--AI collaborative development model.
We demonstrate the system's scalability across three heterogeneous
exemplar projects---a gradient descent optimization study (58 tests,
96.6\% coverage), a categorical linguistics manuscript with 25+
DisCoPy-generated figures (257 tests, \textasciitilde96\% coverage), and
this self-referential architectural analysis (36 tests)---achieving
100\% pipeline success rates. The Standalone Project Paradigm enables
researchers to add arbitrarily many independent studies without
modifying the infrastructure layer, while the Thin Orchestrator pattern
ensures that all domain logic resides in testable \texttt{src/} modules
and scripts serve only as stateless wiring between infrastructure
services and project-specific code.
\end{frame}

\end{document}
